%% Exam randomization
% Examples in RandIntTest.tex
\usepackage{tikz,pgf}
\usepackage{etoolbox,xparse}
\usetikzlibrary{math}
\usepackage{xintfrac}

% MS Co-Pilot and Github Co-Pilot were used to help generate this code.

\newcommand{\randIntLower}{1}   %Lower bound on rand ints
\newcommand{\randIntUpper}{5}   %Upper bound on rand ints

%%%%%%%%%%%%%%%%%%%%%%%%%%%%%%%%%%%%%%%%%%%%%%%%%%%%%%%%%%%%%%%%%%%%%%%%%%%%%%%%
% TODO Move into math121_sp26_exam01_ch1_KEY.tex only?
% Resample random numbers #1 and #2 until distinct
\newcommand*{\resampleUntilnzSum}[2]{%
  \pgfmathtruncatemacro{\again}{#1+#2==0}%   test if #1+#2==0
  \whiledo{\again=1}{%
    \randInts[nz]{#1,#2}% re-gen #1 and #2
    \pgfmathtruncatemacro{\again}{#1+#2==0}% test if #1+#2==0
  }%
}

%%%%%%%%%%%%%%%%%%%%%%%%%%%%%%%%%%%%%%%%%%%%%%%%%%%%%%%%%%%%%%%%%%%%%%%%%%%%%%%%
% \ExplSyntaxOn
% \NewDocumentCommand \defVar { m m }{
%   \fp_set:Nn #1 { \fp_eval:n #2 }
% }
\newcommand{\defVar}[2]{%
%TODO switch to expl3 version after math 121 exam 2 completed
  \pgfmathparse{#2}%
  \pgfmathsetmacro{#1}{\pgfmathresult}%
}

% \ExplSyntaxOn
% \NewDocumentCommand \makeFrac{ m m }{
%% Make a reduced fraction with denominator in [num+1,2*num-1]
%% (vice-versa for num negative)
%   \randInts[nz]{#1}% numerator
%   %TODO #1 will not refer to the cs_name as before
%   \randInts[nz]< \fp_eval:n { #1+abs(#1)/#1 } >( \fp_eval:n { 2*#1-abs(#1)/#1 ){#2}
%   \defVar{\gcdValue}{gcd(#1,#2)}%
%   \defVar{#1}{int(#1/\gcdValue)}%
%   \defVar{#2}{int(#2/\gcdValue)}%
% }
%%%%%%%%%%%%%%%%%%%%%%%%%%%%%%%%%%%%%%%%%%%%%%%%%%%%%%%%%%%%%%%%%%%%%%%%%%%%%%%%
\newcommand{\makeFrac}[2]{%
  %TODO
  \randInts[nz]{#1}% numerator
  % Pick denominator from [num+1,2*num-1] when num>0
  % or [2*num+1,num-1] when num<0
  {% Limit scope of \randIntLower/Upper
    \pgfmathsetmacro{\randIntLower}{\fpeval{#1+abs(#1)/#1}}%
    \pgfmathsetmacro{\randIntUpper}{\fpeval{2*#1-abs(#1)/#1}}%
    \randInts[nz]{#2}% denominator
  }
  \defVar{\gcdValue}{gcd(#1,#2)}%
  \defVar{#1}{int(#1/\gcdValue)}%
  \defVar{#2}{int(#2/\gcdValue)}%
}


\ExplSyntaxOn
%%%%%%%%%%%%%%%%%%%%%%%%%%%%%%%%%%%%%%%%%%%%%%%%%%%%%%%%%%%%%%%%%%%%%%%%%%%%%%%%
% Print negative numbers inside parenthesis
\NewDocumentCommand \printWithParens { m }{
  \fp_compare:nNnTF {#1} < {0}{
    \ensuremath{(#1)}
  }{
    \ensuremath{#1}
  }
}

%%%%%%%% \newcommand{\printWithParens}[1]{%
%%%%%%%%   \pgfmathparse{#1}%
%%%%%%%%   \ifdim\pgfmathresult pt<0pt%
%%%%%%%%     (\pgfmathprintnumber{#1})%
%%%%%%%%   \else%
%%%%%%%%     \pgfmathprintnumber{#1}%
%%%%%%%%   \fi%
%%%%%%%% }
%%%%%%%%%%%%%%%%%%%%%%%%%%%%%%%%%%%%%%%%%%%%%%%%%%%%%%%%%%%%%%%%%%%%%%%%%%%%%%%%
%TODO switch to expl3 version after math 121 exam 2 completed
\newcommand{\randIntTemp}{}     %Used for defining values in randInts
\newcommand{\randInts}[2][]{%
  % #1 = option (default: []; Use nz for nonzero)
  % #2 = list of variables
  % Uses \randIntLower and \randIntUpper as bounds (change with renewcommand if needed)
  \foreach \var in {#2}{%
    \ifstrequal{#1}{nz}{%
      % nonzero method: produces +/- nonzero random integers
      \pgfmathparse{int(ifthenelse(rand > 0, 1, -1)*random(\randIntLower,\randIntUpper))}%
      \expandafter\xdef\var{\pgfmathresult}%
    }{%
      \pgfmathrandominteger{\randIntTemp}{\randIntLower}{\randIntUpper}%
      \expandafter\xdef\var{\randIntTemp}%
    }%
  }%
}
%%%%%%%%%%%%%%%%%%%%%%%%%%%%%%%%%%%%%%%%%%%%%%%%%%%%%%%%%%%%%%%%%%%%%%%%%%%%%%%%
%TODO switch to this version after math 121 exam 2 completed
%%%%%%%% \NewDocumentCommand \randInts { s o D<>{\randIntLower} D(){\randIntUpper} m O{1} }{
%%%%%%%% % \randInts[nz?]<min>(max){name1,name2,...}[t]
%%%%%%%% % - [nz]: produce pm(magnitude), where magnitude in [min,max]
%%%%%%%% % - <min> and (max) are *separate* optional delimited args; defaults are 1 and 5
%%%%%%%% % - clist of names; defines \name globally to the drawn integer
%%%%%%%% % - optional randInt multiplier t (defaults to 1)
%%%%%%%%
%%%%%%%%   \tl_clear_new:N \l__my_number_tl
%%%%%%%%   \seq_clear_new:N \l_numbers_seq
%%%%%%%%   \tl_if_eq:nnT {#2} {nz}{
%%%%%%%%     \int_step_inline:nnn {-#4}{-#3}
%%%%%%%%       { \seq_put_right:Nn \l_numbers_seq { \fp_eval:n{ ##1 * #6 } } }
%%%%%%%%     }
%%%%%%%%   \int_step_inline:nnn {#3}{#4}
%%%%%%%%     { \seq_put_right:Nn \l_numbers_seq { \fp_eval:n{ ##1 * #6 } } }
%%%%%%%%   \seq_shuffle:N \l_numbers_seq
%%%%%%%%
%%%%%%%%   \clist_map_inline:nn { #5 }
%%%%%%%%   {
%%%%%%%%     \IfBooleanTF{#1}{
%%%%%%%%       \seq_pop_left:NN \l_numbers_seq \l__my_number_tl
%%%%%%%%       \cs_gset:cpx { ##1 } { \tl_use:N \l__my_number_tl }
%%%%%%%%     }{
%%%%%%%%       \cs_gset:cpx { ## 1 } { \seq_rand_item:N \l_numbers_seq }
%%%%%%%%     }
%%%%%%%%   }
%%%%%%%% }
%%%%%%%%%%%%%%%%%%%%%%%%%%%%%%%%%%%%%%%%%%%%%%%%%%%%%%%%%%%%%%%%%%%%%%%%%%%%%%%%
\NewDocumentCommand \printIfTrue { m m }{
  \ifnum #1 #2 \fi
  }
%%%%%%%%%%%%%%%%%%%%%%%%%%%%%%%%%%%%%%%%%%%%%%%%%%%%%%%%%%%%%%%%%%%%%%%%%%%%%%%%
\NewDocumentCommand \printIfFalse { m m }{
  \ifnum #1 \else #2 \fi
  }
%%%%%%%%%%%%%%%%%%%%%%%%%%%%%%%%%%%%%%%%%%%%%%%%%%%%%%%%%%%%%%%%%%%%%%%%%%%%%%%%
\NewDocumentCommand \lcm { m m m }{
  \defVar{#1}{ #2*#3/gcd(#2,#3)}%
}
%%%%%%%%%%%%%%%%%%%%%%%%%%%%%%%%%%%%%%%%%%%%%%%%%%%%%%%%%%%%%%%%%%%%%%%%%%%%%%%%
% Translates 1,...,26 to corresponding capital letter
\NewDocumentCommand \numToLetter { m }{
  \pgfmathparse{int(65 + #1 - 1)}\expandafter\char\pgfmathresult
}
%%%%%%%%%%%%%%%%%%%%%%%%%%%%%%%%%%%%%%%%%%%%%%%%%%%%%%%%%%%%%%%%%%%%%%%%%%%%%%%%
% Translates letters to numbers
\NewDocumentCommand \letterToNum { m m }{
  \pgfmathparse{int("`#2"-64)}
  \pgfmathsetmacro{#1}{\pgfmathresult}
}
%%%%%%%%%%%%%%%%%%%%%%%%%%%%%%%%%%%%%%%%%%%%%%%%%%%%%%%%%%%%%%%%%%%%%%%%%%%%%%%%
\bool_new:N \l_exam_print_sign_bool
\bool_new:N \l_exam_signonly_units_bool
% \printCoeff * [signs] {<coeff>}
%   star: suppress leading '+'
%   [signs]: for ±1, print just the sign; otherwise normal numbers
\NewDocumentCommand \printCoeff { s O{} m }{
  \exam_print_coeff:nnn {#1} {#2} {#3}
}

\cs_new_protected:Npn \exam_print_coeff:nnn #1 #2 #3
  {
    % Evaluate coefficient once via pgf
    % TODO Replace pgfmathparse with xparse
    % \pgfmathparse{#3}
    % \tl_set:Nx \l_tmpa_tl { \pgfmathresult } % textual form for pgf printer
    \tl_set:Nx \l_tmpa_tl { \fp_eval:n { #3 } } % textual form for pgf printer
    \fp_set:Nn \l_tmpa_fp { \l_tmpa_tl }        % numeric form for comparisons

    % Star: default signed; star => suppress leading '+'
    \bool_set_true:N \l_exam_print_sign_bool
    \bool_if:nT {#1} { \bool_set_false:N \l_exam_print_sign_bool }

    % Optional keyword: [signs] (case-insensitive; spaces ignored if present)
    \bool_set_false:N \l_exam_signonly_units_bool
    \str_if_eq:eeTF
      { \tl_lower_case:n { \tl_trim_spaces:n {#2} } }
      { signs }
      { \bool_set_true:N \l_exam_signonly_units_bool }
      { }

    % 0 -> print nothing
    \fp_compare:nTF { \l_tmpa_fp == 0 }
      { \ensuremath{} }
      {
        % -1
        \fp_compare:nTF { \l_tmpa_fp == -1 }
          {
            \bool_if:NTF \l_exam_signonly_units_bool {
              \ensuremath{-} } { \ensuremath{-1} }
          }
          {
            % +1
            \fp_compare:nTF { \l_tmpa_fp == 1 }
              {
                \bool_if:NTF \l_exam_signonly_units_bool
                  {  \bool_if:NTF \l_exam_print_sign_bool { % '+' or empty
                    \ensuremath{+} } { \ensuremath{} }
                  }{ \bool_if:NTF \l_exam_print_sign_bool {
                    \ensuremath{+1} } { \ensuremath{1} } }
              }{
                % general numeric case
                \group_begin:
                  %TODO \fp_format:nn
                  \bool_if:NTF \l_exam_print_sign_bool
                    { \ensuremath{\pgfmathprintnumber[print~sign]{\l_tmpa_tl}} }
                    { \ensuremath{\pgfmathprintnumber{\l_tmpa_tl}} }
                \group_end:
              }
          }
      }
  }
%%%%%%%%%%%%%%%%%%%%%%%%%%%%%%%%%%%%%%%%%%%%%%%%%%%%%%%%%%%%%%%%%%%%%%%%%%%%%%%%
\NewDocumentCommand{\signOnly}{m}{%
  \printCoeff[signs]{#1/abs(#1)}
}
%%%%%%%%%%%%%%%%%%%%%%%%%%%%%%%%%%%%%%%%%%%%%%%%%%%%%%%%%%%%%%%%%%%%%%%%%%%%%%%%
\NewExpandableDocumentCommand{\printFrac}{O{}mm}{%
  {%
    %#1: Optional string to trigger display style
    %#2: numerator
    %#3: denominator
    %# Pretty prints the fraction #2/#3
    \str_if_eq:nnT{#1}{d}{\displaystyle}
    \xintTeXsignedFrac{%
      \xintPIrr{\fpeval{#2}/\fpeval{#3}}
    }
  }
}
%%%%%%%%%%%%%%%%%%%%%%%%%%%%%%%%%%%%%%%%%%%%%%%%%%%%%%%%%%%%%%%%%%%%%%%%%%%%%%%%
% Error message for unsupported relations
\msg_new:nnn {exam/signed-ineq}{unknown-relation}
  { The~relation~'#1'~is~not~supported.~Use~<,>,\string\leq,\string\geq. }

% Internal: throw the error
\cs_new_protected:Npn \exam_signed_ineq_error:n #1
  { \msg_error:nnn {exam/signed-ineq}{unknown-relation}{#1} }

% Internal: swap the relation (flip inequality)
\cs_new:Npn \exam_signed_ineq_swap:n #1
  {
    \tl_if_eq:nnTF {#1}{\leq}{\geq}{
    \tl_if_eq:nnTF {#1}{\geq}{\leq}{
    \tl_if_eq:nnTF {#1}{<}{>}{
    \tl_if_eq:nnTF {#1}{>}{<}{
      \exam_signed_ineq_error:n {#1}
    }}}}
  }

% Core: choose swapped or original relation based on the sign of the number
% Uses floating-point comparison so it works for integers and reals.
\cs_new_protected:Npn \exam_signed_ineq:nn #1#2
  {
    \fp_compare:nTF { (#1) < 0 }
      { \exam_signed_ineq_swap:n {#2} }
      { #2 }
  }

% Swap #1 (\geq,\leq,<,>) based on sign of #1
% \signedIneq{<number>}{<relation>}
\NewDocumentCommand{\signedIneq}{ m m }
  { \exam_signed_ineq:nn {\fpeval{#1}}{#2} }
%%%%%%%%%%%%%%%%%%%%%%%%%%%%%%%%%%%%%%%%%%%%%%%%%%%%%%%%%%%%%%%%%%%%%%%%%%%%%%%%

% Scratch variables
\int_zero_new:N \l__bin_len_int
\int_zero_new:N \l__bin_mod_int
\int_zero_new:N \l__bin_pad_int
\tl_new:N        \l__bin_tl

% Pad a *binary string* to the next multiple of 4 (nibble)
\cs_new:Npn \bin_pad_to_nibble:n #1
  {
    % IMPORTANT: expand the input (#1) so macros like \pgfmathresult are resolved
    \tl_set:Nx \l__bin_tl { #1 }

    % Sanitize: trim edges and remove internal spaces
    \tl_trim_spaces:N \l__bin_tl
    \tl_remove_all:Nn \l__bin_tl { \c_space_token }

    % Count and compute padding
    \int_set:Nn \l__bin_len_int { \tl_count:N \l__bin_tl }
    \int_set:Nn \l__bin_mod_int { \int_mod:nn { \l__bin_len_int } { 4 } }

    \int_compare:nNnT { \l__bin_mod_int } > { 0 }
      {
        \int_set:Nn \l__bin_pad_int { \int_eval:n { 4 - \l__bin_mod_int } }
        \int_step_inline:nn { \l__bin_pad_int }
          { \tl_put_left:Nn \l__bin_tl { 0 } }
      }

    % Output
    \tl_use:N \l__bin_tl
  }

% Decimal -> binary -> pad
\cs_new:Npn \bin_of_num_pad_to_nibble:n #1
  {
    % \pgfmathparse{bin(#1)}%
    % \bin_pad_to_nibble:n { \pgfmathresult }
    \bin_pad_to_nibble:n { \int_from_bin:n { #1 } }
  }

% User-facing wrappers
\NewDocumentCommand \binpadnibble { m }        { \bin_pad_to_nibble:n { #1 } }
\NewDocumentCommand \binpadnibblefromnum { m } { \bin_of_num_pad_to_nibble:n { #1 } }
%%%%%%%%%%%%%%%%%%%%%%%%%%%%%%%%%%%%%%%%%%%%%%%%%%%%%%%%%%%%%%%%%%%%%%%%%%%%%%%%

% Internal state (unchanged)
\bool_new:N \l_exam_first_bool
\int_new:N  \l_exam_i_int
\int_new:N  \l_exam_start_int
\int_new:N  \l_exam_stop_int
\int_new:N  \l_exam_step_int

% 1) Numeric range: \myRange{start}{stop}{step}
\NewDocumentCommand{\myRange}{mmm}
 { \exam_myrange_numbers_nomathrel:nnn {#1}{#2}{#3} }

\cs_new_protected:Npn \exam_myrange_numbers_nomathrel:nnn #1#2#3
 {
   \int_set:Nn \l_exam_start_int {#1}
   \int_set:Nn \l_exam_stop_int  {#2}
   \int_set:Nn \l_exam_step_int  {#3}
   \bool_set_true:N \l_exam_first_bool

   % If step = 0, do nothing
   \int_if_zero:nTF { \l_exam_step_int } { }
     {
       % Let the stepper handle direction and stopping
       \int_step_inline:nnnn { \l_exam_start_int } { \l_exam_step_int } { \l_exam_stop_int }
         {
           \bool_if:NF \l_exam_first_bool { , }
           \int_to_arabic:n { ##1 }
           \bool_set_false:N \l_exam_first_bool
         }
     }
 }

% 2) Capital letters: \myLetters[<step>]{<StartLetter>}{<StopLetter>}
\NewDocumentCommand{\myLetters}{O{1}mm}
 { \exam_myrange_letters_nomathrel:nnn {#1}{#2}{#3} }

\cs_new_protected:Npn \exam_myrange_letters_nomathrel:nnn #1#2#3
 {
   \int_set:Nn \l_exam_step_int  {#1}
   \int_set:Nn \l_exam_start_int {`#2}  % ASCII code
   \int_set:Nn \l_exam_stop_int  {`#3}
   \bool_set_true:N \l_exam_first_bool

   \int_if_zero:nTF { \l_exam_step_int } { }
     {
       \int_step_inline:nnnn { \l_exam_start_int } { \l_exam_step_int } { \l_exam_stop_int }
         {
           \bool_if:NF \l_exam_first_bool { , }
           \char_generate:nn { ##1 } {12}
           \bool_set_false:N \l_exam_first_bool
         }
     }
 }

\NewDocumentCommand{\myLettersNum}{O{1}mm}
 {
   \exam_myrange_letters_num:nnn {#1}{#2}{#3}
 }

\cs_new_protected:Npn \exam_myrange_letters_num:nnn #1#2#3
 {
   \int_set:Nn \l_exam_step_int  {#1}
   % Convert alphabetical index to ASCII code: 1→65, 2→66, ...
   \int_set:Nn \l_exam_start_int { 64 + #2 }
   \int_set:Nn \l_exam_stop_int  { 64 + #3 }
   \bool_set_true:N \l_exam_first_bool

   % Do nothing if step = 0
   \int_if_zero:nTF { \l_exam_step_int } { }
     {
       \int_step_inline:nnnn { \l_exam_start_int } { \l_exam_step_int } { \l_exam_stop_int }
         {
           \bool_if:NF \l_exam_first_bool { , }
           \char_generate:nn { ##1 } {12}
           \bool_set_false:N \l_exam_first_bool
         }
     }
 }

%%%%%%%%%%%%%%%%%%%%%%%%%%%%%%%%%%%%%%%%%%%%%%%%%%%%%%%%%%%%%%%%%%%%%%%%%%%%%%%%
% Set random seed based on year, semester, and version
% e.g. VD in spring 2026 corresponds to 1326
% year:           26 : \int_mod:nn { \year }{ 100 }
% exam version   300 : 100 * ( \verNum - 1 )
% semester      1000 : \int_eval:n { \month * 100 + \day

\int_new:N \l__exam_year_two_int
\int_new:N \l__exam_md_int
\int_new:N \l__exam_sem_offset_int
\int_new:N \l__exam_seed_int

% two-digit year
\int_set:Nn \l__exam_year_two_int { \int_mod:nn { \year } { 100 } }

% month-day as MMDD integer (e.g. May 14 -> 514)
\int_set:Nn \l__exam_md_int { \int_eval:n { \month * 100 + \day } }

% semester-based offset: before May 15 -> 1000, before Jun 25 -> 2000,
% before Aug 15 -> 3000, otherwise 4000
\int_set:Nn \l__exam_sem_offset_int { 4000 }
\int_compare:nNnTF { \l__exam_md_int } < { 515 }
  { \int_set:Nn \l__exam_sem_offset_int { 1000 } }          %spring
  {
    \int_compare:nNnTF { \l__exam_md_int } < { 625 }
      { \int_set:Nn \l__exam_sem_offset_int { 2000 } }      %summer I
      {
        \int_compare:nNnTF { \l__exam_md_int } < { 815 }
          { \int_set:Nn \l__exam_sem_offset_int { 3000 } }  %summer II
          { \int_set:Nn \l__exam_sem_offset_int { 4000 } }  %fall
      }
  }

% base seed = semester offset + two-digit year
\int_set:Nn \l__exam_seed_int { \l__exam_sem_offset_int + \l__exam_year_two_int }

% if a `\version` macro exists, convert to numeric and offset by 100*(ver-1)
\cs_if_exist:NTF \version
  {
    \letterToNum{\verNum}{\version}
    \int_set:Nn \l__exam_seed_int { \int_eval:n { \l__exam_seed_int + 100 * ( \verNum - 1 ) } }
  }
  { }

% Set global random seed (query with \sys_rand_seed:
\sys_gset_rand_seed:n { \int_use:N \l__exam_seed_int }
%%%%%%%%%%%%%%%%%%%%%%%%%%%%%%%%%%%%%%%%%%%%%%%%%%%%%%%%%%%%%%%%%%%%%%%%%%%%%%%%
%%%%%%%%%%%%%%%%%%%%%%%%%%%%%%%%%%%%%%%%%%%%%%%%%%%%%%%%%%%%%%%%%%%%%%%%%%%%%%%%
%%%%%%%%%%%%%%%%%%%%%%%%%%%%%%%%%%%%%%%%%%%%%%%%%%%%%%%%%%%%%%%%%%%%%%%%%%%%%%%%
%legacy seed setting
\pgfmathsetseed{ \int_use:N \l__exam_seed_int }
%%%%%%%%%%%%%%%%%%%%%%%%%%%%%%%%%%%%%%%%%%%%%%%%%%%%%%%%%%%%%%%%%%%%%%%%%%%%%%%%
%%%%%%%%%%%%%%%%%%%%%%%%%%%%%%%%%%%%%%%%%%%%%%%%%%%%%%%%%%%%%%%%%%%%%%%%%%%%%%%%
%%%%%%%%%%%%%%%%%%%%%%%%%%%%%%%%%%%%%%%%%%%%%%%%%%%%%%%%%%%%%%%%%%%%%%%%%%%%%%%%
\ExplSyntaxOff%


%%%%%%%%%%%%%%%%%%%%%%%%%%%%%%%%%%%%%%%%%%%%%%%%%%%%%%%%%%%%%%%%%%%%%%%%%%%%%%%%
%%%%%%%%%%%%%%%%%%%%%%%%%%%%%%% old seed setting %%%%%%%%%%%%%%%%%%%%%%%%%%%%%%%
%%%%%%%%%%%%%%%%%%%%%%%%%%%%%%%%%%%%%%%%%%%%%%%%%%%%%%%%%%%%%%%%%%%%%%%%%%%%%%%%
%%%%%%%%%%%%%%%%%%%%%%%%%% % Swap \pdfrandomseed with a number for consistency between compiles
%%%%%%%%%%%%%%%%%%%%%%%%%% \pgfmathsetseed{\number\pdfrandomseed} %\pdfrandomseed
%%%%%%%%%%%%%%%%%%%%%%%%%
%%%%%%%%%%%%%%%%%%%%%%%%%% % Set random seed based on year and semester (eg sp26 is 1026)
%%%%%%%%%%%%%%%%%%%%%%%%%% \defVar{\examRandSeed}{mod(\the\year,100)+int(ifthenelse(\month<6,1,2)*1000)}
%%%%%%%%%%%%%%%%%%%%%%%%%
%%%%%%%%%%%%%%%%%%%%%%%%%% \ExplSyntaxOn
%%%%%%%%%%%%%%%%%%%%%%%%%% \cs_new:Npn \exam_version_to_int:n #1
%%%%%%%%%%%%%%%%%%%%%%%%%%   {
%%%%%%%%%%%%%%%%%%%%%%%%%%     \str_case:nnF {#1}
%%%%%%%%%%%%%%%%%%%%%%%%%%       { {A}{0} {B}{100} {C}{200} {D}{300} }
%%%%%%%%%%%%%%%%%%%%%%%%%%       { \msg_error:nnn { exam } { unknown-version } {#1} 0 }
%%%%%%%%%%%%%%%%%%%%%%%%%%   }
%%%%%%%%%%%%%%%%%%%%%%%%%% \cs_generate_variant:Nn \exam_version_to_int:n { e } % <-- creates :e
%%%%%%%%%%%%%%%%%%%%%%%%%
%%%%%%%%%%%%%%%%%%%%%%%%%% \cs_if_exist:NTF {\version}{ %Modify examRandSeed if \version defined
%%%%%%%%%%%%%%%%%%%%%%%%%%   % \textbackslash examRandSeed=\pgfmathprintnumber{\fpeval{\examRandSeed+\exam_version_to_int:e { \version } } }
%%%%%%%%%%%%%%%%%%%%%%%%%%   \defVar{\examRandSeed}{\number\examRandSeed+\exam_version_to_int:e { \version } }
%%%%%%%%%%%%%%%%%%%%%%%%%% }{}
%%%%%%%%%%%%%%%%%%%%%%%%%% \ExplSyntaxOff
%%%%%%%%%%%%%%%%%%%%%%%%%
%%%%%%%%%%%%%%%%%%%%%%%%%% \pgfmathsetseed{\number\examRandSeed}